\documentclass{amsart}
% For dealing with references we use the comment environment
\usepackage{verbatim}
\newenvironment{reference}{\comment}{\endcomment}
%\newenvironment{reference}{}{}
\newenvironment{slogan}{\comment}{\endcomment}
\newenvironment{history}{\comment}{\endcomment}

% For commutative diagrams we use Xy-pic
\usepackage[all]{xy}

% We use 2cell for 2-commutative diagrams.
\xyoption{2cell}
\UseAllTwocells

% We use multicol for the list of chapters between chapters
\usepackage{multicol}

% This is generally recommended for better output
\usepackage{lmodern}
\usepackage[T1]{fontenc}

% For cross-file-references

% Package for hypertext links:
\usepackage{hyperref}

% For any local file, say "hello.tex" you want to link to please

% Theorem environments.
%
\theoremstyle{plain}
\newtheorem{theorem}[subsection]{Theorem}
\newtheorem{proposition}[subsection]{Proposition}
\newtheorem{lemma}[subsection]{Lemma}

\theoremstyle{definition}
\newtheorem{definition}[subsection]{Definition}
\newtheorem{example}[subsection]{Example}
\newtheorem{exercise}[subsection]{Exercise}
\newtheorem{situation}[subsection]{Situation}

\theoremstyle{remark}
\newtheorem{remark}[subsection]{Remark}
\newtheorem{remarks}[subsection]{Remarks}

\numberwithin{equation}{subsection}

% Macros
%
\def\lim{\mathop{\mathrm{lim}}\nolimits}
\def\colim{\mathop{\mathrm{colim}}\nolimits}
\def\Spec{\mathop{\mathrm{Spec}}}
\def\Hom{\mathop{\mathrm{Hom}}\nolimits}
\def\Ext{\mathop{\mathrm{Ext}}\nolimits}
\def\SheafHom{\mathop{\mathcal{H}\!\mathit{om}}\nolimits}
\def\SheafExt{\mathop{\mathcal{E}\!\mathit{xt}}\nolimits}
\def\Sch{\mathit{Sch}}
\def\Mor{\mathop{\mathrm{Mor}}\nolimits}
\def\Ob{\mathop{\mathrm{Ob}}\nolimits}
\def\Sh{\mathop{\mathit{Sh}}\nolimits}
\def\NL{\mathop{N\!L}\nolimits}
\def\CH{\mathop{\mathrm{CH}}\nolimits}
\def\proetale{{pro\text{-}\acute{e}tale}}
\def\etale{{\acute{e}tale}}
\def\QCoh{\mathit{QCoh}}
\def\Ker{\mathop{\mathrm{Ker}}}
\def\Im{\mathop{\mathrm{Im}}}
\def\Coker{\mathop{\mathrm{Coker}}}
\def\Coim{\mathop{\mathrm{Coim}}}

% Boxtimes
%
\DeclareMathSymbol{\boxtimes}{\mathbin}{AMSa}{"02}

%
% Macros for moduli stacks/spaces
%
\def\QCohstack{\mathcal{QC}\!\mathit{oh}}
\def\Cohstack{\mathcal{C}\!\mathit{oh}}
\def\Spacesstack{\mathcal{S}\!\mathit{paces}}
\def\Quotfunctor{\mathrm{Quot}}
\def\Hilbfunctor{\mathrm{Hilb}}
\def\Curvesstack{\mathcal{C}\!\mathit{urves}}
\def\Polarizedstack{\mathcal{P}\!\mathit{olarized}}
\def\Complexesstack{\mathcal{C}\!\mathit{omplexes}}
% \Pic is the operator that assigns to X its picard group, usage \Pic(X)
% \Picardstack_{X/B} denotes the Picard stack of X over B
% \Picardfunctor_{X/B} denotes the Picard functor of X over B
\def\Pic{\mathop{\mathrm{Pic}}\nolimits}
\def\Picardstack{\mathcal{P}\!\mathit{ic}}
\def\Picardfunctor{\mathrm{Pic}}
\def\Deformationcategory{\mathcal{D}\!\mathit{ef}}

\setlength{\emergencystretch}{3em}
\begin{document}
\title[Stacks Project, Tag 0ECN]{309 --- Finiteness of the next Serre-depth defect locus}
\maketitle
\noindent Copyright (C) 2005--2025 Johan de Jong. Stacks Project authors. Source: \href{https://stacks.math.columbia.edu/tag/0ECN}{Tag 0ECN}.

\noindent Editorial difficulty note (not author text). Difficulty H2: The proof reduces the next-depth defect to the preceding Serre level along a carefully chosen regular element. It then chooses that element so its complement lies entirely in the Cohen-Macaulay locus, leaving only the finitely many induction defects. This is a global finiteness argument for a depth-stratified locus..

\noindent Editorial provenance note (not author text). History: excerpt from revision \texttt{a0444\allowbreak 6e57e\allowbreak c1fbc\allowbreak 252a8\allowbreak 71afc\allowbreak ec775\allowbreak 2fb28\allowbreak 07b14}, local-cohomology.tex, lines 1904--1965. Reference presentation was adapted; original-author context and core support blocks were retained or added. The mathematical wording is unchanged. Licensed under GNU FDL 1.2 or later; no invariant sections or cover texts. See ../licenses/COPYING. Context blocks reproduce author definitions and supporting results; they are not additional counted problems.

\begin{definition}
\label{algebra-definition-conditions}
Let $R$ be a Noetherian ring.
Let $k \geq 0$ be an integer.
\begin{enumerate}
\item We say $R$ has property {\it $(R_k)$} if for every prime $\mathfrak p$
of height $\leq k$ the local ring $R_{\mathfrak p}$ is regular.
We also say that $R$ is {\it regular in codimension $\leq k$}.
\item We say $R$ has property {\it $(S_k)$} if for every prime $\mathfrak p$
the local ring $R_{\mathfrak p}$ has depth at least
$\min\{k, \dim(R_{\mathfrak p})\}$.
\item Let $M$ be a finite $R$-module. We say $M$ has property $(S_k)$
if for every prime $\mathfrak p$ the module
$M_{\mathfrak p}$ has depth at least
$\min\{k, \dim(\text{Supp}(M_{\mathfrak p}))\}$.
\end{enumerate}
\end{definition}

\begin{lemma}
\label{lemma-finite-nr-points-next-S}
Let $X$ be a Noetherian scheme with dualizing complex $\omega_X^\bullet$.
Let $\mathcal{F}$ be a coherent $\mathcal{O}_X$-module. Let $k \geq 0$
be an integer. Assume $\mathcal{F}$ is $(S_k)$.
Then there is a finite number of points $x \in X$ such that
$$
\text{depth}(\mathcal{F}_x) = k
\quad\text{and}\quad
\dim(\text{Supp}(\mathcal{F}_x)) > k
$$
\end{lemma}

\begin{proof}
We will prove this lemma by induction on $k$. The base case $k = 0$
says that $\mathcal{F}$ has a finite number of embedded associated points,
which follows from Divisors, Lemma \href{https://stacks.math.columbia.edu/tag/05AF}{lemma finite ass (Tag 05AF)}.

\medskip\noindent
Assume $k > 0$ and the result holds for all smaller $k$.
We can cover $X$ by finitely many affine opens, hence we may
assume $X = \Spec(A)$ is affine. Then $\mathcal{F}$ is the
coherent $\mathcal{O}_X$-module associated to a finite $A$-module $M$
which satisfies $(S_k)$. We will use
Algebra, Lemmas \href{https://stacks.math.columbia.edu/tag/0B52}{lemma one equation module (Tag 0B52)} and
\href{https://stacks.math.columbia.edu/tag/090R}{lemma depth drops by one (Tag 090R)}
without further mention.

\medskip\noindent
Let $f \in A$ be a nonzerodivisor on $M$. Then $M/fM$ has $(S_{k - 1})$.
By induction we see that there are finitely many
primes $\mathfrak p \in V(f)$ with
$\text{depth}((M/fM)_\mathfrak p) = k - 1$ and
$\dim(\text{Supp}((M/fM)_\mathfrak p)) > k - 1$.
These are exactly the primes $\mathfrak p \in V(f)$ with
$\text{depth}(M_\mathfrak p) = k$ and
$\dim(\text{Supp}(M_\mathfrak p)) > k$.
Thus we may replace $A$ by $A_f$ and $M$ by $M_f$
in trying to prove the finiteness statement.

\medskip\noindent
Since $M$ satisfies $(S_k)$ and $k > 0$ we see that $M$ has no
embedded associated primes
(Algebra, Lemma \href{https://stacks.math.columbia.edu/tag/031Q}{lemma criterion no embedded primes (Tag 031Q)}).
Thus $\text{Ass}(M)$ is the set of generic points of the support
of $M$. Thus Dualizing Complexes, Lemma \href{https://stacks.math.columbia.edu/tag/0EHS}{lemma CM open (Tag 0EHS)}
shows the set
$U = \{\mathfrak q \mid M_\mathfrak q\text{ is Cohen-Macaulay}\}$
is an open containing $\text{Ass}(M)$.
By prime avoidance (Algebra, Lemma \href{https://stacks.math.columbia.edu/tag/00DS}{lemma silly (Tag 00DS)})
we can pick $f \in A$ with
$f \not \in \mathfrak p$ for $\mathfrak p \in \text{Ass}(M)$
such that $D(f) \subset U$.
Then $f$ is a nonzerodivisor on $M$
(Algebra, Lemma \href{https://stacks.math.columbia.edu/tag/00LD}{lemma ass zero divisors (Tag 00LD)}).
After replacing $A$ by $A_f$ and $M$ by $M_f$ (see above) we
find that $M$ is Cohen-Macaulay.
Thus for all $\mathfrak q \subset A$ we have
$\dim(M_\mathfrak q) = \text{depth}(M_\mathfrak q)$
and hence the set described in the lemma is empty
and a fortiori finite.
\end{proof}
\end{document}
